\Artigo*%
{historia}%
{\textit{Wow!}}%
{Santiago González Gómez}%
{De como un código alfanumérico impreso polo radiotelescopio Big Ear segue
fascinándonos case 50 anos despois.}%

\begin{multicols}{2}

15 de agosto, 1977. Jerry Ehman, que traballaba como voluntario no
radiotelescopio Big Ear da Universidade Estatal de Ohio, revisaba as lecturas
do aparello. O Big Ear adicábase a escanear o ceo, recibir sinais de radio, e %adicar: forma menos recomendábel por dedicar
imprimir as lecturas en follas de papel. Máis concretamente, este
radiotelescopio buscaba sinais de banda estreita, que aínda que poden ocorrer
na natureza, son emisións que adoitan asociarse con seres intelixentes. Os
datos que revisaba Ehman referíanse á intensidade de sinal recibida por
diversas canles, que se imprimía como un ratio sinal/fondo. Así, por exemplo,
se se rexistraba un sinal 4 veces máis forte que o ruído de fondo, imprimíase
un catro. Para aforrar espazo, os números do 10 en adiante representábanse con
letras: A para 10, B para 11, e así sucesivamente. De súpeto, Ehman topouse coa
seguinte lectura, agora mundialmente famosa:

\vspace*{-0.8em}

\begin{center}
    6 E Q U J 5.
\end{center}

\vspace*{-0.8em}

Falamos dun sinal que, no seu pico, chegou a $30{,}5 \pm 0{,}5$ a intensidade do
fondo. Ehman rodeouno e comentouno cunha única palabra que acabou por darlle
nome: \textit{Wow!}.

\begin{figure}[H]
    \centering
    \includegraphics[width=\linewidth]{./revistas/005/imaxes/wow.png}
    \caption{
        Detalle do printout no que Ehman observou o sinal
        \textit{Wow!} \cite{Ehman}.
    }
    \label{im:wow}
\end{figure}

\subsection*{O sinal}

Pero, que \textit{escoitou} o Big Ear naquel verán de 1977? Para podermos
comezar coas conxecturas, cómpre coñecermos algo máis do sinal. Na figura
\ref{im:perfil} podemos observar unha recreación da forma do sinal a partir
dos puntos medidos. O sinal durou uns 72 segundos (pasaron uns 12 segundos
entre cada par de medicións), e a súa forma encaixa co obxectivo do
radiotelescopio pasando por diante dunha fonte fixa no ceo. Aínda que a gráfica
resulta lixeiramente asimétrica, isto explícase polo deseño do Big Ear e nada
ten que ver co sinal orixinal \cite{Kraus}. A súa frecuencia é duns $1420$
MHz, terriblemente preto da frecuencia da transición hiperfina do hidróxeno.
Ademais, o sinal \textit{Wow!} está moi localizado no espectro (menos de $10$
kHz de largo de banda).

$1420$ MHz é unha frecuencia de grande importancia en astronomía, tanto que
está prohibido empregala na Terra para transmisións de radio para non interferir
coas medicións radioastronómicas. Primeiro, porque se observa moito: o
hidróxeno é o elemento máis común no Universo e ademais pode emitir esta
radiación espontaneamente. Segundo, porque corresponde a unha onda de radio que
pode atravesar masas de po opacas á luz, permitíndonos ver o que hai detrás.

\begin{figure}[H]
    \centering
    \includegraphics[width=0.9\linewidth]{./revistas/005/imaxes/perfil.png}
    \caption{Perfil do sinal. (Foto: Wikipedia)}
    \label{im:perfil}
\end{figure}

Precisamente que o sinal tivese esta frecuencia reforza a hipótese de atribuílo
a unha raza extraterrestre. Os seres humanos sabemos da importancia desta
frecuencia como padrón natural do cosmos, e cando precisamos dunha unidade de
tempo o máis universal posible coa que encriptar as mensaxes que mandamos nas
nosas sondas espaciais por se as atopasen intelixencias extraterrestres (en
particular, as placas da Pioneer e o disco de ouro da Voyager), elixiuse
precisamente a inversa desta frecuencia como unidade. Non sería esperable que
seres intelixentes doutros planetas emitisen os seus sinais intergalácticos
nesta frecuencia por idénticas razóns?\footnote{Nun artigo 18 anos anterior á
    detección do sinal \textit{Wow!}, proponse $1420$ MHz como a frecuencia
máis prometedora para procurar mensaxes dalgunha sociedade extraterrestre
\cite{morrison1959searching}.}

Se nos inclinamos pola teoría de que o sinal \textit{Wow!} ten orixe
extraterrestre, unha pregunta natural sería, codificaba algún tipo de
información? A resposta é que, se o facía, non chegou ata nós. Para obter
información dunha onda, esperariamos que esta estivese modulada, que oscilase
entre un par de frecuencias ou amplitudes, e isto non é o que escoitou o Big
Ear. Porén, isto non quere dicir que o sinal non estivese modulado: lembremos
que o radiotelescopio medía cada 12 segundos, e a información púidose perder
nesa fiestra temporal. Precisamente a «falta de información» codificada no
sinal \textit{Wow!} é o que deixa aberta a posibilidade moi real de que en
realidade poida ser explicada cunha orixe natural.

Outra pregunta que aínda non respondemos sobre o sinal é tamén moi natural:
ben, de onde procedía? Lamentablemente, é outra da que non sabemos a resposta,
alomenos con gran precisión. Resulta que o Big Ear estaba deseñado de tal xeito %alomenos: forma menos recomendábel por "polo menos"
que escaneaba dúas rexións celestiais simultaneamente, e non temos maneira de
distinguir en cal das dúas rexións se atopaba a fonte do sinal. Iso si, este
deseño permítenos facer unha reflexión curiosa: resulta que o movemento do Big
Ear era tal que a segunda rexión volvía pasar por onde estivera a primeira 3
minutos antes. Como o sinal \textit{Wow!} só se detectou unha vez, temos unha
de dúas posibilidades: ou a fonte transmitiu durante un período indeterminado
antes do escáner da primeira rexión e extinguiuse nos 3 minutos que tardou en
«chegar» a segunda rexión, ou a transmisión iniciouse nese intervalo de 3
minutos entre escaneamentos \cite{Lemmino}. Sexa como for, todo apunta a que o
sinal \textit{Wow!} foi un evento moi puntual, que tivemos moita sorte de
captar e, en efecto, durante case 50 anos fomos incapaces de recibir outro sinal
similar destas dúas rexións.

\subsection*{As posibles orixes}

Imaxinemos que, como bos escépticos, non nos fiamos de que o sinal
\textit{Wow!} sexa de orixe extraterrestre. Que outras posibles explicacións
temos? Ben, dende logo, non esperamos que a fonte sexa unha estrela ou un
planeta, que debería de emitir en continuo (polo tanto, nun amplo largo de
banda). Tamén foron propostos e descartados un par de cometas como
responsables.

Como a frecuencia da liña hiperfina do hidróxeno está protexida na Terra,
tampouco esperamos que nós sexamos os causantes. Con todo, non corramos tanto a
descartar esta hipótese. Que esta frecuencia estea internacionalmente prohibida
non quere dicir que ninguén a puidese usar. Porén, outros problemas descartan a
posibilidade dunha orixe humana. Ningún satélite nin aeronave coñecida se
atopaban nesa zona do ceo cando se mediu o sinal \textit{Wow!}, e aínda que o
estivesen, calquera destes dous obxectos moveríase demasiado rápido para
podermos recrear o padrón medido. Tampouco é realista que o que se medise fose
un sinal emitido no chan que se reflectiu nalgún anaco de lixo espacial,
basicamente polas características dinámicas que tería que ter a órbita deste
espello para observar o sinal da figura \ref{im:perfil}.

Unha teoría amplamente discutida é a do escintileo interestelar. Puidera ser
que o sinal fose amplificado por algún efecto de tipo coherente ao atravesar o
medio interestelar dende a súa orixe ata nós. Porén, o problema desta e outras
teorías é que non resolven a cuestión da orixe. A relevancia da teoría do
escintileo estelar é que permitiría a transformación dun sinal continuo de
maior largo de banda no sinal observado, que por suposto é un mecanismo que non
pode ignorarse.

Se cadra a explicación natural máis satisfactoria é a de
\cite{mendez2024arecibo}. Atopando similitudes entre o sinal \textit{Wow!} e
radiación moito menos intensa proveniente de nubes de hidróxeno, os seus
autores teorizan que o sinal \textit{Wow!} puido ser causado pola reemisión,
por parte dunha destas nubes de hidróxeno, dalgún fenómeno superradiante, como
un chorro dun magnétar (un tipo de púlsar cun fortísimo campo magnético). %non estou seguro de se habería que dicir magnétar ou magnetar, déixoo así por se acaso

\begin{figure}[H]
    \centering
    \includegraphics[width=0.75\linewidth]{./revistas/005/imaxes/reemision.png}
    \caption{
        Fenómeno de amplificación e reemisión dun chorro radiante por parte dunha
        nube de hidróxeno \cite{mendez2024arecibo}.
    }
    \label{im:reemision}
\end{figure}

O certo é que, a un ano do 50º aniversario do sinal \textit{Wow!}, seguimos sen
ter respostas concretas para moitas das súas incógnitas. Aínda que unha orixe
extraterrestre non está descartada, en palabras de Ehman, non podemos «draw
vast conclusions from half-vast data». Sen máis deteccións similares, só queda
especular...

\printbibliography

\end{multicols}
